\section*{ВВЕДЕНИЕ}
\addcontentsline{toc}{section}{ВВЕДЕНИЕ}

 Одним из основных инструментов обеспечения конфиденциальности данных является криптография. Современные стандарты шифрования позволяют защитить информацию от несанкционированного доступа как при её передаче по открытым каналам связи, так и при хранении на носителях.

Среди существующих алгоритмов симметричного шифрования наиболее распространенным и надежным признан стандарт AES (Advanced Encryption Standard). Он является обязательным для использования в государственных структурах многих стран и широко применяется в коммерческом секторе для защиты персональных данных, финансовых транзакций и интеллектуальной собственности.

Использование языка программирования Common Lisp для реализации криптографических инструментов позволяет задействовать преимущества функциональной парадигмы программирования: высокую степень абстракции, модульность и удобство работы со сложными структурами данных. Разработка собственной программной библиотеки, реализующей стандарт AES-128, является актуальной задачей для понимания механизмов защиты информации и создания гибких инструментов безопасности.

\emph{Цель настоящей работы} – разработка программной библиотеки на языке Common Lisp, реализующей алгоритм AES-128 для защиты информации путем шифрования данных с возможностью их последующего дешифрования. 

Для достижения поставленной цели необходимо решить \emph{следующие задачи:}
\begin{itemize}
	\item изучить теоретические основы и математическую модель алгоритма AES-128;
	\item спроектировать модульную архитектуру библиотеки в среде Common Lisp;
	\item реализовать алгоритм расширения ключа и раундовые преобразования данных;
	\item разработать механизмы обработки данных произвольной длины с использованием стандарта PKCS\#7;
	\item произвести тестирование корректности работы библиотеки.
\end{itemize}

\emph{Структура и объем работы.} Отчет состоит из введения, 4 разделов основной части, заключения, списка использованных источников, 2 приложений. 

\emph{Во введении} сформулирована цель работы, поставлены задачи разработки, описана актуальность и структура работы.

\emph{В первом разделе} приводится теоретическое описание алгоритма AES-128, рассматриваются этапы преобразования данных и математический аппарат шифрования.

\emph{Во втором разделе} на стадии технического проектирования описывается архитектура библиотеки, структура программных модулей и выбор инструментальных средств разработки.

\emph{В третьем разделе} приводится детальное описание реализации ключевых модулей библиотеки: генерации ключей, функций замен и перестановок, а также логики дополнения блоков.

\emph{В четвертом разделе} описывается процесс проведения тестов, приводятся контрольные примеры работы библиотеки и результаты верификации алгоритма.

В заключении излагаются основные результаты работы, полученные в ходе разработки.

В приложении А представлен графический материал (схемы алгоритма).